\subsection{The all-Vd0 kernel}

\begin{definition}[V-gaps and active-gap rank]
\label{def:body-active-gap-rank}
On the boundary edge $e_{i,i+1}$, let $\Gamma_i$ be the closed set left
uncovered by the two incident open vertex roles.  Equivalently, with the
standard edge parameter,
\[
 \Gamma_i=
 \begin{cases}
  [B_i,1-A_{i+1}],&B_i\le1-A_{i+1},\\
  \varnothing,&B_i>1-A_{i+1}.
 \end{cases}
\]
A nonempty $\Gamma_i$ is a \emph{V-gap}; equality gives a singleton V-gap.
In a hypothetical cover every V-gap lies in a center trace.  The
\emph{active-gap rank} $\mathrm{gr}$ is the number of positive center traces
that contain a V-gap.  Hence
\[
 \mathrm{gr}\in\{0,1,2\},
 \qquad
 \mathrm{gr}\le1\quad\hbox{in CE1}.
\]
\end{definition}

The all-Vd0 cases are organized by one small table.
\[
\begin{array}{c|ccc}
 &\mathrm{gr}=0&\mathrm{gr}=1&\mathrm{gr}=2\\ \hline
N_+=0&\text{strict cycle}&\text{one-gap endpoints}&
       \text{paired endpoints}\\
N_+=1&\text{nine-point obstruction}&\text{five-V-triangle return}&
       \text{paired endpoints}.
\end{array}
\]
The last column is CE2-only.  Thus Strategy~2 does not need a different
calculus for each cell: it uses one transfer interface with three terminal
shapes.

For $N_+=0$ and $\mathrm{gr}=0$, the vertex traces alone cover every boundary
edge.  Nonsupercriticality gives $1-B_i\ge A_i$, while open boundary coverage
gives $A_{i+1}>1-B_i$.  Hence
\[
 A_{i+1}>1-B_i\ge A_i,
\]
and the strict cycle is immediate.  No transfer map is needed in this cell.  For $N_+=1$ and $\mathrm{gr}=0$, no
transfer calculation is needed; this cell is owned by Strategy~4.

